%---------------------------------------------------------
\section{Background}
\label{sec:background}

Functional Data Analysis (FDA) treats each observed trajectory---rather than
each individual measurement---as the fundamental unit of statistical
analysis. A discretely sampled time series is first converted into a
continuous function by expressing it as a linear combination of basis
elements, and all subsequent analysis (smoothing, principal component
analysis, regression) is carried out on these functional objects
\cite{ramsay2005}. This basis-expansion step is not a neutral
preprocessing choice: the basis system determines which features of the
underlying process are preserved, which are treated as noise and
suppressed, and which patterns become visible to downstream methods such as
Functional Principal Component Analysis (FPCA).

Most applications of FDA to macroeconomic time series adopt smooth,
globally supported bases---typically B-splines or truncated Fourier
series---implicitly assuming that the underlying trajectories belong to a
Sobolev space of functions with bounded, uniformly distributed roughness
\cite{ramsay2005}. This approach is applied by \cite{padilla2020} to
annual GDP growth rates of Latin American economies, representing each country's
growth trajectory with a B-spline basis before conducting FPCA to study
cross-country modes of variation. This choice is reasonable for economies
whose growth process evolves smoothly, but Latin American GDP growth
series are also punctuated by abrupt, localized shocks---the 1982 debt
crisis, the 1999 regional currency crises, the 2008 global financial
crisis, and the 2020 COVID-19 contraction---whose magnitude and timing
differ sharply across countries.

An alternative to globally smooth bases is provided by wavelets, which
form an unconditional basis for Besov spaces
$B^{s}_{p,q}$---function classes that, unlike Sobolev spaces, tolerate
spatially inhomogeneous regularity, i.e.\ functions that are smooth over
most of their domain but exhibit sharp, localized transitions elsewhere
\cite{donoho1998}. Because wavelets localize simultaneously in time and
frequency, they can in principle represent a smooth trend and a sharp
crisis episode within the same trajectory without forcing a single
global smoothness parameter to compromise between the two. This
theoretical property motivates a systematic empirical comparison between
globally smooth bases (B-splines, Fourier) and a multiresolution basis
(wavelets) for representing volatile economic series.

\section{Problem statement}
\label{sec:problem_statement}

The choice of functional basis in FDA applications to economic time
series is rarely treated as an object of empirical scrutiny in its own
right: a single basis system---most often B-splines---is selected on
convenience or precedent, fitted, and used as the input to FPCA or other
downstream analyses, without assessing whether its underlying smoothness
assumptions are appropriate for the data at hand. This is problematic for
series such as GDP growth in volatile emerging economies, where
country-level heterogeneity in volatility (from Guatemala's comparatively
stable trajectory to Venezuela's crisis-dominated one) and time-varying,
crisis-driven variance call into question the adequacy of a single global
smoothness criterion applied uniformly across countries.

Two consequences follow when this choice is left unexamined. First, a
basis that over-smooths may achieve a favorable reconstruction error while
silently discarding the crisis episodes that are often the most
economically meaningful features of the series, so that low error is
mistaken for a faithful representation. Second, because FPCA is computed
directly on the fitted functional objects, any distortion introduced at
the basis-expansion stage propagates into the estimated modes of
variation, so that the resulting principal components may reflect
artifacts of the smoothing choice rather than genuine cross-country
structure. Existing work extending the study of \cite{padilla2020} has not
systematically compared alternative basis families under a common,
fair evaluation protocol that jointly considers reconstruction fidelity,
coefficient economy (sparsity), and the interpretability of the resulting
FPCA. This thesis addresses that gap by comparing B-spline, Fourier, and
wavelet basis expansions on an extended panel of Latin American GDP growth
series (1960--2024).

\section{Objectives}
\label{sec:objectives}

\subsection{General Objective}

To comparatively evaluate B-spline, Fourier, and wavelet basis expansions
for representing the GDP growth trajectories of volatile Latin American
economies, in terms of reconstruction fidelity, coefficient economy, and
their effect on the interpretability of the resulting functional
principal component analysis.

\subsection{Specific Objectives}

\begin{enumerate}
    \item To extend the annual GDP growth dataset of \cite{padilla2020}
    for 20 Latin American countries through 2024, incorporating the
    COVID-19 contraction and subsequent recovery period.

    \item To implement three families of functional basis expansions
    for the resulting panel---B-spline, Fourier, and wavelet (Haar, DB2,
    Sym4)---under a common, reproducible computational pipeline.

    \item To define and compute comparison metrics---reconstruction
    error (RMSE), sparsity, and FPCA variance explained---that allow the
    three basis families to be evaluated at matched levels of
    reconstruction quality.

    \item To assess, for each basis family, the extent to which crisis
    episodes and country-specific volatility structure are preserved or
    suppressed in the functional reconstruction.

    \item To characterize how the choice of basis affects the parsimony
    and economic interpretability of the resulting FPCA modes of
    variation, relative to the original results of \cite{padilla2020}.
\end{enumerate}

%---------------------------------------------------------

%\section{Justification}
